
\subsubsection{Findings}

The mechanism validation components provide evidence that compositional tokenization can preserve architecture-specific signals (per-family $\rho \geq$ 0.68 when trained on single-family subsets) while maintaining family structure in learned embeddings (silhouette score 0.52). The approach showed statistically significant improvements over flat-weight concatenation ($\Delta\rho$=0.18, p=0.0005) and engineered statistical features ($\Delta\rho$=0.12, p=0.008).

However, overall performance on the heterogeneous test set ($\rho$=0.294) was substantially below the target threshold. The gap between per-family performance ($\rho \geq$ 0.68) and overall heterogeneous performance ($\rho$=0.294) suggests that cross-architecture transfer at the tested 150-model scale may be limited.

\subsubsection{Limitations}

\textbf{Scale}: The experiment used 150 models (150 train, 30 test) rather than the initially planned 750 models (600 train, 150 test). The small test set (30 samples) contributes to wide confidence intervals and limited statistical power. Per-family ablation results ($\rho \geq$ 0.68) suggest that mechanism validation is successful within families, but overall performance indicates that larger training populations may be necessary for robust cross-architecture learning.

\textbf{Domain Shift}: Generalization gaps were computed on CIFAR-10 for models pretrained on ImageNet (CNNs and ViTs) or trained on MNIST (MLPs), due to ImageNet dataset unavailability. This domain mismatch may affect measured generalization characteristics and confound quality signals.

\textbf{Transformer Tokenization}: Vision Transformer processing ($\rho$=0.68) showed slightly lower performance than CNNs ($\rho$=0.72) and MLPs ($\rho$=0.75) in per-family ablation. Q/K/V matrix extraction may not fully capture multi-head attention structure, or the domain shift effect may differentially impact Vision Transformers.

\textbf{Scope}: Validation focused on image classification models (CNNs, Vision Transformers, MLPs). Generalization to other domains (NLP transformers, reinforcement learning policies) or other model properties (robustness, calibration) remains untested.

\subsubsection{Comparison to Related Work}

The compositional approach differs from prior work by addressing heterogeneous model zoos rather than homogeneous collections (NFT on MNIST MLPs) or single architecture families (DWSNets for CNNs). The mechanism validation provides proof-of-concept evidence at 150-model scale, though performance levels remain below those reported in prior work on homogeneous settings.
